\section{Appendix: API and Data Flow}
\subsection{Main API routes}
\begin{table}[ht]
\centering\small
\begin{tabularx}{\textwidth}{@{}l l X@{}}
\toprule
\textbf{Method} & \textbf{Path} & \textbf{Purpose}\\
\midrule
GET & \texttt{/api/stats}, \texttt{/api/daily} & Return catalogue counts and the deterministic daily problem.\\
GET & \texttt{/api/collections} & Search (accent-insensitive) and filter collection records by kind, grade, and page.\\
GET & \texttt{/api/collections/\{id\}} & Return one collection with ordered problems and source metadata.\\
GET & \texttt{/api/problems} & Search individual questions by text, grade, canonical topic, and page.\\
GET & \texttt{/api/lessons} & Return or search lesson records (with examples and Markdown body).\\
GET & \texttt{/api/books}, \texttt{/api/docs} & Return or search books, magazines, and PDF documents.\\
POST & \texttt{/api/register}, \texttt{/api/login} & Create account or start a signed HttpOnly cookie session.\\
POST & \texttt{/api/logout}; GET \texttt{/api/me} & End the session or read the signed-in account.\\
POST & \texttt{/api/change-password} & Change password for the signed-in account.\\
POST & \texttt{/api/request-reset}, \texttt{/api/reset-password} & Generate a 30-minute reset token and set a new password.\\
GET/POST & \texttt{/api/photos} & List photos (optionally filtered by \texttt{problem}) or upload compressed photos.\\
GET/DELETE & \texttt{/api/photos/\{id\}} & View (\texttt{nosniff}) or delete an account-owned image.\\
\bottomrule
\end{tabularx}
\caption{Main API routes in the basic project.}
\end{table}

\subsection{Simple request example}
When a student searches for a collection, the browser reads the search text and grade selection. It makes a URL such as \texttt{/api/collections?search=Ha+Noi\&grade=9\&page=1}. FastAPI reads these query parameters, selects matching JSON records, and returns JSON. The browser then creates collection cards from the returned records. In the photo flow, the browser sends file bytes instead of JSON text; the API checks the session and stores the bytes in PostgreSQL.

\subsection{Environment configuration}
The application needs two private environment values: \texttt{DATABASE\_URL}, the connection string from the PostgreSQL provider, and \texttt{SESSION\_SECRET}, a long random signing value. A sample names-only file is provided as \texttt{.env.example}. The real local \texttt{.env.local} file is ignored by Git. Neither private value should be copied into this report, screenshots, or public source control.
